\Qhead{3.1}{Случайные величины и процессы. Усреднение и эргодичность}{Случайные величины и процессы, способы их описания. Стационарный случайный процесс. Статистическое усреднение и усреднение во времени. Эргодичность. Измерение вероятностей и средних значений.}
\Subfile{3.1.1}{sections/03_statistics/q01/s01}
\Subfile{3.1.2}{sections/03_statistics/q01/s02}

\recap{Процесс описывается ансамблем реализаций и многовременными распределениями. Стационарность — инвариантность статистики, эргодичность — возможность заменить ансамбль одной достаточно длинной записью.}
\board{Среднее, ковариация, двойной интеграл дисперсии среднего и контрпример случайной константы.}
\oral{Достаточно ли удвоить частоту оцифровки, чтобы вдвое уменьшить дисперсию среднего при неизменной длительности записи?}
\bait{Можно оставить контрпример эргодичности: случайная постоянная имеет стационарную статистику, но не усредняется за время.}
\sources{Биленко и соавт.~\cite{bilenko}, §§4.1.1--4.1.3; университетский курс математической статистики~\cite{msustats}, определения распределений и оценок. Дисперсия временного среднего выведена здесь через ковариацию.}
